% Source: Wald GR, physical PDF p. 28 (printed p. 15).
% Coverage: Figure 2.3, directional derivatives in a coordinate chart.
% Figures/tables/footnotes: Figure only; no tables or footnotes.
% Status: source-reviewed and compile-clean.
% Uncertainties: none.
\begingroup
\renewcommand{\thefigure}{2.3}
\begin{figure}[htbp]
  \centering
\begin{tikzpicture}[x=0.72cm,y=0.72cm,>=Stealth,font=\small]
  \draw[rounded corners=7pt] (0.2,3.1)
    .. controls (-0.2,4.1) and (0.4,4.9) .. (1.2,4.75)
    .. controls (2.2,5.25) and (2.4,4.3) .. (3.3,4.4)
    .. controls (4.4,4.4) and (4.7,3.6) .. (4.6,2.7)
    .. controls (4.8,1.8) and (3.5,1.3) .. (2.7,1.7)
    .. controls (1.5,1.2) and (0.1,1.8) .. (0.2,3.1);
  \node at (0.75,4.1) {\(M\)};
  \draw (1.65,2.95) ellipse (0.83 and 0.72);
  \fill (2.02,3.47) circle (1.5pt);
  \node[left=1pt] at (2.00,3.47) {\(p\)};
  \node at (1.42,2.86) {\(\mathcal O\)};

  \draw (6.0,2.1) rectangle (9.6,5.0);
  \node at (6.48,4.62) {\(\mathbb{R}^n\)};
  \draw (8.20,3.62) ellipse (0.95 and 0.64);
  \node at (8.05,4.00) {\(U\)};
  \fill (7.90,3.47) circle (1.5pt);
  \node[right=1pt] at (7.94,3.47) {\(\psi(p)\)};
  % Coordinate axes of the chart, drawn clear of \(U\).
  \draw[->] (6.75,2.45) -- (6.75,3.18);
  \draw[->] (6.75,2.45) -- (7.52,2.45);
  \node[left=1pt] at (6.72,3.06) {\(x^2\)};
  \node[below=1pt] at (7.46,2.42) {\(x^1\)};
  \draw[->] (2.70,3.47) -- (7.80,3.47);
  \node[above=1pt] at (5.25,3.50) {\(\psi\)};
  \draw[->] (2.02,3.32) -- (2.02,-1.86);
  \node[right=1pt] at (2.06,0.30) {\(f\)};
  \draw (0.70,-1.95) -- (3.70,-1.95);
  \node[above right=1pt] at (2.55,-1.92) {\(\mathbb{R}\)};
\end{tikzpicture}
  \caption{A diagram illustrating the definition of the directional derivatives,
    \(X_\mu\), used in Theorem~\ref{thm:2.2.1}.}
  \label{fig:2.3}
\end{figure}
\endgroup
